\section{Introduction}\label{sec:introduction}

Multi-turn dialogue is the dominant mode of interaction with large language models (LLMs), and system prompts that specify task goals form the foundation of instruction following. As conversations grow longer, however, models progressively lose track of these goal specifications. \citet{dongre2026attention} quantified this degradation across 12 models: the Goal Attention Ratio (GAR), which measures how much attention the model allocates to goal tokens in the system prompt, declines by an average of 70.8\% between early and mid-conversation turns. We replicate this finding on Qwen-2.5-7B-Instruct, where GAR drops from 0.0418 at turn~1 to 0.0122 at mid-conversation, consistent with the cross-model pattern. This attention decay represents a systematic failure of multi-turn models to maintain task coherence, and it grows worse with each additional turn.

A striking aspect of this degradation is what it does \emph{not} affect. \citet{dongre2026attention} showed that goal-relevant information persists in the residual stream throughout the conversation: linear discriminant analysis probes recover goal content with AUC\,=\,0.98 even at late turns where attention to goal tokens has collapsed. The information is present in the hidden states, but the attention mechanism fails to route it toward the output computation. This dissociation between information availability and information utilization suggests a natural intervention: if the residual stream already carries goal information, a learnable gate on the residual pathway could amplify or attenuate layer contributions to recover goal attention. This reasoning motivates ThreadGate, which applies a per-layer sigmoid gate $g_l$ to each residual update: $x_{l+1} = x_l + g_l \cdot f_l(x_l)$, following the ReZero residual scaling principle \citep{bachlechner2021rezero}.

\begin{figure*}[t]
  \centering
  \includegraphics[width=\textwidth]{fig1_overview.pdf}
  \caption{Overview of ThreadGate and its failure mode. (Left) ThreadGate adds per-layer sigmoid gates to the residual stream, trained jointly with an auxiliary GAR loss that targets goal attention. (Right) Under teacher-forced evaluation, GAR-ours reaches 0.49 across all fine-tuned variants; under autoregressive generation, it drops to 0.016 for ThreadGate, a 30$\times$ gap. The four-way ablation (ThreadGate, Uniform Gate, LoRA-only, Base) shows that gating contributes zero measurable benefit over LoRA-only fine-tuning.}
  \label{fig:overview}
\end{figure*}

ThreadGate combines per-layer gates with a joint training objective: next-token prediction (NTP), an auxiliary GAR loss that pushes attention toward goal tokens via an MSE target ($\tau = 0.5$), and a gate diversity loss that prevents all gates from collapsing to a single value (\S\ref{sec:method}). The final loss formulation emerged from seven design iterations (v1 through v7) spanning four loss families (cumulative, maximize, smooth saturation, and MSE target); only the MSE target with gate diversity converged stably. To isolate the contribution of gating from the effects of fine-tuning and auxiliary losses, we evaluate a four-way ablation: ThreadGate (per-layer gates), Uniform Gate (a single shared gate), LoRA-only (gates fixed at 1.0, same loss terms), and the unmodified Base model. This factorial design supports three targeted comparisons that separate per-layer conditioning, gating as a mechanism, and fine-tuning itself (\S\ref{sec:experiments}).

The training phase succeeds: GAR rises from 0.02 to 0.48 over three epochs. Teacher-forced evaluation on held-out data confirms that the learned attention pattern generalizes within the training distribution, with GAR-ours reaching 0.49 to 0.50 across all fine-tuned variants. The failure emerges under autoregressive generation, where GAR-ours drops to 0.016 for ThreadGate, a 30$\times$ gap from the teacher-forced setting. The four-way ablation reveals that this gap is independent of the gating mechanism: ThreadGate, Uniform Gate, and LoRA-only achieve indistinguishable GAR-ours (spread of 0.75 percentage points) and identical perplexity (2.78) on held-out data. Gating contributes zero measurable benefit. To understand why, we trace a causal chain through the training dynamics (\S\ref{sec:analysis}). Gradient conflict between $\mathcal{L}_\text{NTP}$ and $\mathcal{L}_\text{GAR}$ is pervasive (68 to 79\% of layers show negative cosine similarity), but a PCGrad \citep{yu2020pcgrad} experiment eliminates direction conflict without rescuing GAR convergence, establishing gradient conflict as a symptom rather than a cause. The root failure is exposure bias \citep{bengio2015scheduled}: attention patterns optimized under teacher-forced prefixes do not transfer to model-generated text, and the 2 to 7\% gate attenuation at inference is far too weak to compensate for this distributional mismatch (Figure~\ref{fig:overview}).

Our contributions are as follows:
\begin{enumerate}
    \item We provide the first empirical characterization of how exposure bias destroys the generalization of auxiliary attention losses: teacher-forced GAR and autoregressive GAR differ by a factor of 30$\times$, while the primary NTP objective remains robust. This demonstrates that token-level self-recovery from exposure bias \citep{he2019exposure} does not extend to auxiliary objectives that depend on distributional properties of the input prefix.
    \item We design a PCGrad diagnostic exclusion experiment that establishes a complete diagnostic chain: gradient direction conflict between NTP and GAR losses is a symptom of multi-objective optimization dynamics, not the mechanism responsible for GAR failure (\S\ref{sec:analysis}).
    \item ThreadGate and the systematic four-way ablation provide a principled negative baseline for residual gating in multi-turn attention recovery, with seven loss formulations and three gating configurations ruled out (\S\ref{sec:method}, \S\ref{sec:experiments}).
    \item We are conducting cross-model validation on Llama-3.1-8B-Instruct to verify that the convergence pattern and the training-inference gap generalize across model families (\S\ref{sec:cross-model}).
\end{enumerate}